\begin{lexlean}{ChaossOssProjectViabilityStrategyProofs}
\useglossary{lexlean.std.bool@1.2.0}
\useglossary{lexlean.std.nat@1.2.0}
\useglossary{prism.arch@1.0.0}
\importmodule{ChaossOssProjectViabilityStrategy}
\title{Text application model}
\begin{semanticmodule}
definition{
  name: breaksInv001,
  parameters: [
    request : ChaossOssProjectViabilityStrategy::ChaossOssProjectViabilityStrategyRequest
  ],
  result: Bool,
  body: (((((project request admitted) <=? 5) && ((project request refused) <=? 5)) && (1 <=? (project request refused))) && ((ChaossOssProjectViabilityStrategy::chaossOssProjectViabilityStrategyJudge request) =? (ChaossOssProjectViabilityStrategy::chaossOssProjectViabilityStrategyAnswerComplete)))
}

theorem{
  name: inv001Holds,
  parameters: [
    request : ChaossOssProjectViabilityStrategy::ChaossOssProjectViabilityStrategyRequest
  ],
  statement: ((breaksInv001 request) = false),
  proof: linear_arithmetic{
    definitions: [
      ChaossOssProjectViabilityStrategy::chaossOssProjectViabilityStrategyAnswerComplete,
      ChaossOssProjectViabilityStrategy::chaossOssProjectViabilityStrategyAnswerOpen,
      ChaossOssProjectViabilityStrategy::chaossOssProjectViabilityStrategyAnswerPartial,
      ChaossOssProjectViabilityStrategy::chaossOssProjectViabilityStrategyAnswerUnmeasured,
      ChaossOssProjectViabilityStrategy::chaossOssProjectViabilityStrategyJudge,
      ChaossOssProjectViabilityStrategy::chaossOssProjectViabilityStrategyRule1,
      ChaossOssProjectViabilityStrategy::chaossOssProjectViabilityStrategyRule2,
      ChaossOssProjectViabilityStrategy::chaossOssProjectViabilityStrategyRule3,
      breaksInv001
    ]
  },
  axioms: [Quot.sound, propext]
}

definition{
  name: breaksInv002,
  parameters: [
    request : ChaossOssProjectViabilityStrategy::ChaossOssProjectViabilityStrategyRequest
  ],
  result: Bool,
  body: (((((project request admitted) <=? 5) && ((project request refused) <=? 5)) && ((project request admitted) <? 5)) && ((ChaossOssProjectViabilityStrategy::chaossOssProjectViabilityStrategyJudge request) =? (ChaossOssProjectViabilityStrategy::chaossOssProjectViabilityStrategyAnswerComplete)))
}

theorem{
  name: inv002Holds,
  parameters: [
    request : ChaossOssProjectViabilityStrategy::ChaossOssProjectViabilityStrategyRequest
  ],
  statement: ((breaksInv002 request) = false),
  proof: linear_arithmetic{
    definitions: [
      ChaossOssProjectViabilityStrategy::chaossOssProjectViabilityStrategyAnswerComplete,
      ChaossOssProjectViabilityStrategy::chaossOssProjectViabilityStrategyAnswerOpen,
      ChaossOssProjectViabilityStrategy::chaossOssProjectViabilityStrategyAnswerPartial,
      ChaossOssProjectViabilityStrategy::chaossOssProjectViabilityStrategyAnswerUnmeasured,
      ChaossOssProjectViabilityStrategy::chaossOssProjectViabilityStrategyJudge,
      ChaossOssProjectViabilityStrategy::chaossOssProjectViabilityStrategyRule1,
      ChaossOssProjectViabilityStrategy::chaossOssProjectViabilityStrategyRule2,
      ChaossOssProjectViabilityStrategy::chaossOssProjectViabilityStrategyRule3,
      breaksInv002
    ]
  },
  axioms: [Quot.sound, propext]
}

definition{
  name: breaksInv003,
  parameters: [
    request : ChaossOssProjectViabilityStrategy::ChaossOssProjectViabilityStrategyRequest
  ],
  result: Bool,
  body: ((((((project request admitted) <=? 5) && ((project request refused) <=? 5)) && ((project request refused) =? 0)) && (5 <=? (project request admitted))) && (not ((ChaossOssProjectViabilityStrategy::chaossOssProjectViabilityStrategyJudge request) =? (ChaossOssProjectViabilityStrategy::chaossOssProjectViabilityStrategyAnswerComplete))))
}

theorem{
  name: inv003Holds,
  parameters: [
    request : ChaossOssProjectViabilityStrategy::ChaossOssProjectViabilityStrategyRequest
  ],
  statement: ((breaksInv003 request) = false),
  proof: linear_arithmetic{
    definitions: [
      ChaossOssProjectViabilityStrategy::chaossOssProjectViabilityStrategyAnswerComplete,
      ChaossOssProjectViabilityStrategy::chaossOssProjectViabilityStrategyAnswerOpen,
      ChaossOssProjectViabilityStrategy::chaossOssProjectViabilityStrategyAnswerPartial,
      ChaossOssProjectViabilityStrategy::chaossOssProjectViabilityStrategyAnswerUnmeasured,
      ChaossOssProjectViabilityStrategy::chaossOssProjectViabilityStrategyJudge,
      ChaossOssProjectViabilityStrategy::chaossOssProjectViabilityStrategyRule1,
      ChaossOssProjectViabilityStrategy::chaossOssProjectViabilityStrategyRule2,
      ChaossOssProjectViabilityStrategy::chaossOssProjectViabilityStrategyRule3,
      breaksInv003
    ]
  },
  axioms: [Quot.sound, propext]
}

definition{
  name: breaksInv004,
  parameters: [
    request : ChaossOssProjectViabilityStrategy::ChaossOssProjectViabilityStrategyRequest
  ],
  result: Bool,
  body: ((((((project request admitted) <=? 5) && ((project request refused) <=? 5)) && (1 <=? (project request admitted))) && ((project request refused) =? 0)) && ((ChaossOssProjectViabilityStrategy::chaossOssProjectViabilityStrategyJudge request) =? (ChaossOssProjectViabilityStrategy::chaossOssProjectViabilityStrategyAnswerUnmeasured)))
}

theorem{
  name: inv004Holds,
  parameters: [
    request : ChaossOssProjectViabilityStrategy::ChaossOssProjectViabilityStrategyRequest
  ],
  statement: ((breaksInv004 request) = false),
  proof: linear_arithmetic{
    definitions: [
      ChaossOssProjectViabilityStrategy::chaossOssProjectViabilityStrategyAnswerComplete,
      ChaossOssProjectViabilityStrategy::chaossOssProjectViabilityStrategyAnswerOpen,
      ChaossOssProjectViabilityStrategy::chaossOssProjectViabilityStrategyAnswerPartial,
      ChaossOssProjectViabilityStrategy::chaossOssProjectViabilityStrategyAnswerUnmeasured,
      ChaossOssProjectViabilityStrategy::chaossOssProjectViabilityStrategyJudge,
      ChaossOssProjectViabilityStrategy::chaossOssProjectViabilityStrategyRule1,
      ChaossOssProjectViabilityStrategy::chaossOssProjectViabilityStrategyRule2,
      ChaossOssProjectViabilityStrategy::chaossOssProjectViabilityStrategyRule3,
      breaksInv004
    ]
  },
  axioms: [Quot.sound, propext]
}
\end{semanticmodule}
\end{lexlean}
